\chapter{Criterio de continuidad, reapertura y nueva fase del sistema}

\section{Continuidad disciplinada}
\begin{definition}
Se llama continuidad disciplinada del sistema al vínculo explícito mediante el cual MADELEINE conserva memoria editorial, jerarquía funcional y trazabilidad suficiente al pasar desde la clausura editorial provisional de la fase actual hacia reaperturas, reintegraciones o nuevas fases posteriores.
\end{definition}

\begin{observation}
La continuidad no equivale a inmovilidad ni a repetición. Su función es impedir que el crecimiento posterior entre como ruptura silenciosa o como acumulación amorfa.
\end{observation}

\section{Reapertura disciplinada}
\begin{definition}
Existe reapertura disciplinada cuando puede declararse con claridad:
\begin{enumerate}
  \item qué zona se reabre;
  \item por qué se reabre;
  \item qué se conserva intacto;
  \item bajo qué criterio podrá reintegrarse lo revisado.
\end{enumerate}
\end{definition}

\begin{observation}
La reapertura no destruye por sí sola el cierre ya alcanzado. Lo modifica localmente bajo reglas explícitas.
\end{observation}

\section{Nueva fase editorial}
\begin{definition}
Existe nueva fase editorial cuando el sistema puede abrir un nuevo ciclo de expansión sin destruir la arquitectura ya estabilizada entre Tomo I, Tomo II y Tomo III.
\end{definition}

\begin{observation}
Una nueva fase no debe entrar como simple añadido de material. Debe entrar con razón explícita, límites visibles y criterio de incorporación.
\end{observation}

\section{Reintegración del material reabierto}
\begin{definition}
Se llama reintegración disciplinada al retorno explícito de material reabierto al corpus principal, conservando memoria de la intervención realizada y relación legible con la versión previa.
\end{definition}

\begin{observation}
No toda reapertura exige reintegración inmediata. A veces conviene anexo, reserva o desplazamiento a una fase posterior.
\end{observation}

\section{Cláusula de conservación}
\begin{observation}
Toda continuidad, reapertura, nueva fase o reintegración debe conservar, salvo revisión expresamente declarada:
\begin{enumerate}
  \item Tomo I como volumen fundacional principal;
  \item Tomo II como canon intermedio y organización operativa;
  \item Tomo III como capa metodológica ya estabilizada en su función.
\end{enumerate}
\end{observation}

\section{Cláusula de alcance}
\begin{observation}
Este bloque compacta la gramática mínima de continuidad posterior a la clausura editorial provisional. No constituye todavía una teoría exhaustiva de todas las fases futuras posibles del corpus.
\end{observation}

